\begin{frame}{회귀에서 분류로: 같은 도구를 다른 문제에}
  \begin{itemize}
    \item 로지스틱회귀라는 이름 때문에 헷갈리기 쉽지만, 이건 회귀가
      아니라 \textbf{분류} 알고리즘.
    \item 선형회귀의 출력 $\boldsymbol{w}^T\boldsymbol{x}$ 는 $-\infty$ 부터 $+\infty$ 까지
      아무 값이나 되는데, ``스팸일 확률''은 반드시 0과 1 사이여야
      한다.
  \end{itemize}
  \vskip0.8em
  \begin{block}{핵심 아이디어}
    2.1절의 \kb{모델}과 \kb{경사하강법}을 그대로 재사용하고,
    \textbf{``출력을 확률로 바꾸는 변환''}과 \textbf{``그 확률에
    맞는 새 손실함수''}만 추가하면 된다 --- 이번 장 후반부의
    핵심.
  \end{block}
\end{frame}
